\documentclass[11pt,a4paper]{article}
\usepackage[T1]{fontenc}
\usepackage{isabelle,isabellesym}
\usepackage{amsfonts,amsmath,amssymb}

% this should be the last package used
\usepackage{pdfsetup}

% urls in roman style, theory text in math-similar italics
\urlstyle{rm}
\isabellestyle{it}

\begin{document}

\title{The Irrationality of $\zeta(5)$}
\author{Alex Hyman}
\maketitle

\begin{abstract}\sloppy
We formalise the proof that $\zeta(5)$ is irrational from A.~Fauzan's 2026 preprint~\cite{fauzan},
following the Lean~4/Mathlib formalisation by domino14~\cite{lean}, of which this entry is a
port.  The argument builds integer polynomials
$Q_n$ with $0 < Q_n(\zeta(5)) < e^{-cn^2}$ as Hankel determinants of a rational functional: a
positive moment representation of $\zeta(5)$ gives positivity, $p$-adic estimates through
unimodular bases in two prime ranges and a uniform asymptotic for the valuations give the
integrality factor, and a logarithmic-energy estimate for an explicit measure gives the
archimedean bound, proving irrationality.  The main theorem is \isa{zeta 5 $\notin$ $\mathbb{Q}$}
for the \isa{zeta} of the entry Zeta\_\allowbreak Function; the only oracle is \isa{holds}, from the 72 numeric
inequalities discharged by \isa{approximation}.  The port is not literal: Lean's dependently typed
\isa{Fin\ n} matrices and $\Sigma$-indexed bases become Jordan\_\allowbreak Normal\_\allowbreak Form
matrices over explicit index lists, and the $p$-adic field $\mathbb{Q}_p$ is replaced by truncated
rational series with congruences modulo powers of $p$, so that the whole development lives in HOL
over $\mathbb{Q}$ without Types-To-Sets; Lean's kernel computation on 1052 dyadic intervals becomes
\isa{approximation} with splitting.

Thirteen theories are independent of the proof and reusable: a calculus for lower bounds on
$p$-adic valuations of rationals and polynomials with a mod-$p$ unimodularity criterion
(\isa{unimodular\_\allowbreak criterion}), the Vandermonde determinant for Jordan\_\allowbreak Normal\_\allowbreak Form matrices
(\isa{det\_\allowbreak vandermonde\_\allowbreak jnf}) and positivity of Gram determinants, Riemann-sum limits
(\isa{outer\_\allowbreak prime\_\allowbreak sum\_\allowbreak bound}) of sums over primes from the prime number theorem, the
log-sine integral and the circle average of $\ln|z-a|$ (\isa{ln\_\allowbreak sin\_\allowbreak integral},
\isa{circle\_\allowbreak avg\_\allowbreak log}), the arcsine law
and its potential (\isa{arcsine\_\allowbreak potential}), Andr\'eief's identity for $\sigma$-finite measures
(\isa{andreief}), and the zero-mass logarithmic energy inequality on $\mathbb{C}$
(\isa{log\_\allowbreak energy\_\allowbreak zero\_\allowbreak mass\_\allowbreak nonpos}).

The port was an experiment in AI proof engineering: an orchestration model (Claude Fable 5.1)
wrote Isar skeletons with fact hints, proof agents (Claude Opus 5.5) closed about 1,100 leaves
interactively in Isabelle/jEdit, and automated scripts harvested \isa{try0} and
\isa{sledgehammer} results to reduce context load on the agents.  The problem was chosen because
its final statement is short and is expressed with existing AFP definitions, so the result is
easy to audit.  The formalisation effort spanned 3 days of continuous work with a fair amount of
human intervention to keep Fable and the agents on track, give them Isabelle-related hints,
insist on legible proofs, and oversee the structure of the project.  The author has audited the
files and takes full responsibility for their contents.
\end{abstract}




\section*{Relation to the Lean formalisation}

The definitions, statements and proof structure follow the Lean project theory by theory (each
Isabelle theory names its Lean counterparts).  The main departures:
\begin{itemize}\sloppy
\item The $p$-adic field $\mathbb{Q}_p$ of the paper's Sections 3.1--3.2 is not constructed.
  Every $p$-adically convergent series is replaced by its length-$M$ partial sum (the \isa{\_\allowbreak tr}
  definitions) and the identities become congruences modulo powers of $p$, so the development
  stays in HOL over $\mathbb{Q}$ without Types-To-Sets (theories Zeta5\_\allowbreak Truncated\_\allowbreak Expansions,
  Zeta5\_\allowbreak Pole\_\allowbreak Distribution).
\item Lean's \isa{Fin n} matrices and $\Sigma$-indexed bases become Jordan\_\allowbreak Normal\_\allowbreak Form matrices
  with explicit index lists (\isa{inner\_\allowbreak idx}, \isa{outer\_\allowbreak idx}); dependent types were never
  needed, only bookkeeping.
\item Lean checks the potential inequality (A.9) by kernel computation on 1052 dyadic intervals;
  here it is 52 region lemmas by \isa{approximation} with splitting, and the 16 energy enclosures
  (A.10) are per-term \isa{approximation} runs after rescaling the logarithms into $[1,2)$.
\item Mathlib facts with no HOL/AFP counterpart were proved from scratch and are the general
  theories listed in the abstract; the moment representation (Proposition 2.2) is derived from
  the Laplace transform of $\sin$ and the cotangent partial-fraction formula of the AFP rather
  than from a Hurwitz-zeta integral, the prime number theorem is the AFP's
  (\isa{\isasymtheta\_\allowbreak asymptotics}), the von Staudt--Clausen facts come from the entry Bernoulli, and the rational
  valuations from Kummer\_\allowbreak Congruence.
\item The zero-mass energy inequality (Lemma 6.2) is stated with explicit hypotheses (a finite
  dominating measure with $\ln|z-w|$ integrable on the product and a null diagonal), which is
  what the paper's ``finite energy, compact support'' has to mean for signed measures.
\end{itemize}

\tableofcontents
\newpage
\parindent 0pt\parskip 0.5ex

\input{session}

\bibliographystyle{abbrv}
\bibliography{root}

\end{document}
